\documentclass[main.tex]{subfiles}


\begin{document}

% \AddToShipoutPictureFG{%
% \begin{tikzpicture}[overlay,remember picture]
% \path (current page.south west) -- (current page.north east)
% node[midway,scale=1,color=gray,sloped,opacity=0.4] {\textnormal{Кравченко Игорь Игоревич\hspace{1cm} +79010144910\hspace{1cm} super.antig@yandex.ru}};
% \end{tikzpicture}
% }


% %from pagenumber to up
% \phantomsection 
% \hypertarget{toc}{}

 

 

 
   

\section{Система отсчета. Действие}

Прежде всего стоит вспомнить, что \emph{системой отсчета} (СО), попросту говоря, называется тело, к которому прикреплен \emph{наблюдатель}. Так, на рис.\,\ref{fig:p26} показаны наблюдатели, связанные с разными~СО; рядом с ними происходит падение тяжелого шара (положение шара фиксируют ежесекундно).

\begin{figure}[H]
    \centering

    \begin{tikzpicture}[font=\small,            line cap=round,                             >={Stealth[inset=0pt,round,length=3mm, angle'=20]},vec/.style={>={Stealth[inset=0pt,round,length=3.5mm, angle'=20]}}]


\filldraw[lightgray,semitransparent] (0,0) rectangle (12,-0.3);
\draw (0,0) -- (12,0);

\begin{scope}[xshift=8cm]
    
%%motion1
\filldraw[lightgray,draw=black] (0,0.5) rectangle (4,1);
\node [right,black] at (4,.75) {Т};
\draw[->,NavyBlue,thick,vec] (1,1.2) --++ (-1,0) node[black,above] {$\vec v_\text{т}$};
% PERSON
\def\W{5.2}     % ground width
\def\L{3.2}     % length
\def\T{0}    % plank thickness
\def\H{2.0}     % human height
  \def\x{2} % human position
  \coordinate (O) at (0,0);
\begin{scope}[yshift=1cm]  
  \draw[thick] (\x,\H) circle(0.3) coordinate (H) node[xshift=0.3cm,right] {$\text{Н}_2$};
  \draw[thick] (H)++(-90:0.3) coordinate (N) to[out=-92,in=92]++ (0,-0.40*\H) coordinate (P);
  \draw[thick] (N)++(-95:0.03) to[out=-105,in=90]++ (-0.02*\W,-0.4*\H);
  \draw[thick] (N)++(-85:0.03) to[out=-80,in=90]++ (0.02*\W,-0.4*\H);
  \draw[thick] (P) to[out=-92,in=82] (\x-0.02*\W,\T);
  \draw[thick] (P) to[out=-80,in=90] (\x+0.02*\W,\T);
\end{scope}

%wheels
\draw[black,fill=black]
    (1,.3) circle(.3cm)
    (3,.3) circle(.3cm);
\draw[fill=gray]
    (1,.3) circle(.2cm)
    (3,.3) circle(.2cm);
\end{scope}

%%%%%%%%%%%%%%%%%%%%%%%%%%%%%%%%%%
%ball
\shade[ball color=lightgray,nearly transparent,yshift=4cm] (6,.3) circle (.3);
\shade[ball color=lightgray,semitransparent,yshift=3cm] (6,.3) circle (.3);
\shade[ball color=lightgray] (6,.4) circle (.3) node[xshift=.3cm,right] {Ш};

%vec
\node at (6,3.3) [circle,fill=NavyBlue,inner sep=0pt,minimum size=1mm,label=below:] {};
\draw[->,NavyBlue,thick,vec] (6,3.3) --++ (0,-.75) node[black,right] {$\vec v_1$};

\node at (6,0.4) [circle,fill=NavyBlue,inner sep=0pt,minimum size=1mm,label=below:] {};
\draw[->,NavyBlue,thick,vec] (6,0.4) --++ (0,-1.5) node[black,right,inner ysep=0mm] {$\vec v_2$};

%%%%%%%%%%%%%%%%%%%%%%%%%%%%%%%%%%%%
%person 2
% PERSON
\def\W{5.2}     % ground width
\def\L{3.2}     % length
\def\T{0}    % plank thickness
\def\H{2.0}     % human height
  \def\x{3} % human position
  \coordinate (O) at (0,0);
\begin{scope}[yshift=0cm]  
  \draw[thick] (\x,\H) circle(0.3) coordinate (H) node[xshift=0.3cm,right] {$\text{Н}_1$};
  \draw[thick] (H)++(-90:0.3) coordinate (N) to[out=-92,in=92]++ (0,-0.40*\H) coordinate (P);
  \draw[thick] (N)++(-95:0.03) to[out=-105,in=90]++ (-0.02*\W,-0.4*\H);
  \draw[thick] (N)++(-85:0.03) to[out=-80,in=90]++ (0.02*\W,-0.4*\H);
  \draw[thick] (P) to[out=-92,in=82] (\x-0.02*\W,\T);
  \draw[thick] (P) to[out=-80,in=90] (\x+0.02*\W,\T);
\end{scope}



    \end{tikzpicture}
\caption{Наблюдатели и падающий шар}
\label{fig:p26}
\end{figure}

Наблюдатель~$\text{Н}_1$ стоит на поверхности планеты, а наблюдатель~$\text{Н}_2$ стоит на тележке~Т, движущейся с постоянной скоростью~$\vec v_\text{т}$ относительно этой планеты. Тогда можно считать, что имеются две СО: СО~<<планета>> и СО~<<тележка>>\footnote{Это <<правильные>> СО. Дальнейшие рассуждения справедливы для таких СО.}.   

\textbf{Действие}~--- это влияние одного тела, приводящее к \emph{изменению скорости} другого тела.
Можно согласиться, что шар~Ш на рис.\,\ref{fig:p26} главным образом подвергается влиянию планеты, поэтому его скорость и меняется.


Те же наблюдатели рассматривают уже приземлившийся шар (рис.\,\ref{fig:p27}).

\begin{figure}[H]
    \centering

    \begin{tikzpicture}[font=\small,            line cap=round,                             >={Stealth[inset=0pt,round,length=3mm, angle'=20]},vec/.style={>={Stealth[inset=0pt,round,length=3.5mm, angle'=20]}}]


\filldraw[lightgray,semitransparent] (0,0) rectangle (12,-0.3);
\draw (0,0) -- (12,0);

\begin{scope}[xshift=8cm]
%%motion1
\filldraw[lightgray,draw=black] (0,0.5) rectangle (4,1);
\node [right,black] at (4,.75) {Т};
\draw[->,NavyBlue,thick,vec] (1,1.2) --++ (-1,0) node[black,above] {$\vec v_\text{т}$};
% PERSON
\def\W{5.2}     % ground width
\def\L{3.2}     % length
\def\T{0}    % plank thickness
\def\H{2.0}     % human height
  \def\x{2} % human position
  \coordinate (O) at (0,0);
\begin{scope}[yshift=1cm]  
  \draw[thick] (\x,\H) circle(0.3) coordinate (H) node[xshift=0.3cm,right] {$\text{Н}_2$};
  \draw[thick] (H)++(-90:0.3) coordinate (N) to[out=-92,in=92]++ (0,-0.40*\H) coordinate (P);
  \draw[thick] (N)++(-95:0.03) to[out=-105,in=90]++ (-0.02*\W,-0.4*\H);
  \draw[thick] (N)++(-85:0.03) to[out=-80,in=90]++ (0.02*\W,-0.4*\H);
  \draw[thick] (P) to[out=-92,in=82] (\x-0.02*\W,\T);
  \draw[thick] (P) to[out=-80,in=90] (\x+0.02*\W,\T);
\end{scope}

%wheels
\draw[black,fill=black]
    (1,.3) circle(.3cm)
    (3,.3) circle(.3cm);
\draw[fill=gray]
    (1,.3) circle(.2cm)
    (3,.3) circle(.2cm);
\end{scope}

%%%%%%%%%%%%%%%%%%%%%%%%%%%%%%%%%%
%ball
\shade[ball color=lightgray] (6,.3) circle (.3) node[xshift=.3cm,right] {Ш};


%%%%%%%%%%%%%%%%%%%%%%%%%%%%%%%%%%%%
%person 2
% PERSON
\def\W{5.2}     % ground width
\def\L{3.2}     % length
\def\T{0}    % plank thickness
\def\H{2.0}     % human height
  \def\x{3} % human position
  \coordinate (O) at (0,0);
\begin{scope}[yshift=0cm]  
  \draw[thick] (\x,\H) circle(0.3) coordinate (H) node[xshift=0.3cm,right] {$\text{Н}_1$};
  \draw[thick] (H)++(-90:0.3) coordinate (N) to[out=-92,in=92]++ (0,-0.40*\H) coordinate (P);
  \draw[thick] (N)++(-95:0.03) to[out=-105,in=90]++ (-0.02*\W,-0.4*\H);
  \draw[thick] (N)++(-85:0.03) to[out=-80,in=90]++ (0.02*\W,-0.4*\H);
  \draw[thick] (P) to[out=-92,in=82] (\x-0.02*\W,\T);
  \draw[thick] (P) to[out=-80,in=90] (\x+0.02*\W,\T);
\end{scope}



    \end{tikzpicture}
\caption{Наблюдатели и неподвижный шар}
\label{fig:p27}
\end{figure}

На рис.\,\ref{fig:p27} наблюдатель~$\text{Н}_1$ фиксирует постоянную (равную нулю) скорость шара~Ш, а наблюдатель~$\text{Н}_2$~--- постоянную ненулевую. Скорость шара неизменна, хотя и тут \emph{воздействие планеты} на шар нельзя отрицать. Ясно, что имеется еще некоторое действие на шар со стороны \emph{поверхности}. 

Итак, тело может также одновременно испытывать влияния со стороны сразу нескольких других тел. Если же при наличии нескольких действий на тело оно \emph{покоится} или движется \emph{равномерно прямолинейно}, то в этом случае говорят, что эти \textbf{действия скомпенсированы}. 

















 
\end{document}